% Resumo e abstract: versão provisória. Reescrever no fim (S17), quando todos os
% resultados existirem. Limite ABNT: 150 a 500 palavras.

\setlength{\absparsep}{18pt}
\begin{resumo}
Laudos de mamografia apresentam inconsistências documentadas, como baixa adesão ao léxico BI-RADS e grafias divergentes de um mesmo termo, e não há ferramenta que confira automaticamente o que o laudo afirma contra o que a imagem mostra. Este trabalho propõe um sistema híbrido de auditoria que cruza a imagem mamográfica com o laudo já emitido e sinaliza dois tipos de problema: omissões de achados suspeitos e discordâncias de descritores BI-RADS. O sistema combina um detector de lesões treinado no VinDr-Mammo, um classificador de categoria BI-RADS e densidade por mama, um extrator de informação do laudo em português e um auditor baseado em regras rastreáveis. Na validação, o detector de massas alcançou AUC de 0,832 (IC 95\%: 0,768 a 0,889) em nível de mama e encontrou 61,2\% das massas com um falso positivo por imagem; uma ablação de pré-processamento não detectou diferença entre o recorte da mama, a ausência de recorte e o realce por CLAHE. \pendente{completar com os resultados do classificador, da extração dos laudos e da auditoria}

\textbf{Palavras-chave}: mamografia. BI-RADS. auditoria de laudos. visão computacional. detecção de objetos.
\end{resumo}

\begin{resumo}[Abstract]
\begin{otherlanguage*}{english}
\pendente{traduzir o resumo quando a versão final estiver pronta}

\vspace{\onelineskip}
\noindent
\textbf{Keywords}: mammography. BI-RADS. report auditing. computer vision. object detection.
\end{otherlanguage*}
\end{resumo}
